% GENERATED FILE. Edit canonical lessons, then run npm run generate:books.
% canonical-source-hash: fnv1a64:74e94f243b74d73f
% canonical-lessons: RU-C518-ispolnyat, RU-C518-ispolnitel, RU-C518-repetitsiya, RU-C518-repetirovat, RU-C518-gastroli

\chapter{Rehearsals and Tours}
\label{ch:ru-rehearsals-and-tours}
\hlchaptermodality{518}

\begin{chapteropening}
\textbf{By the end of this chapter:} \emph{I can talk about rehearsing and performing.}

Complete the last lesson of chapter 518: I can talk about rehearsing and performing.
\end{chapteropening}

\section[\texorpdfstring{\ru{исполнять} (ispolnyát)}{исполнять (ispolnyát)}]{\ru{исполнять} (ispolnyát) — to perform; to carry out}

\label{lesson:RU-C518-ispolnyat}



\begin{quote}
Before the new one: say the Russian for to conduct, then the Russian for a composer.
\end{quote}

\subsection*{You'll want to know: \ru{исполнять}}
\textbf{\ru{исполнять}} (\emph{ispolnyát}) — \textquotedblleft{}to perform; to carry out\textquotedblright{}.

\begin{quote}
\ru{Певица} \ru{исполняет} \ru{старые} \ru{народные} \ru{песни}.
\end{quote}

\begin{quote}
\emph{Pevítsa ispolnyáyet stáryye naródnyye pésni.}
\end{quote}

\begin{quote}
The singer performs old folk songs.
\end{quote}

Imperfective; perfective \ru{исполнить}. Also \emph{\ru{исполнять} \ru{мечту}} — to make a dream come true.

\begin{cousinweb}
\emph{\ru{исполнять}} is built on \emph{\ru{ис}-} and \emph{\ru{полный}} 'full', so to perform is to fulfil; EtymDB lists the verb without a source.
\end{cousinweb}

\begin{grammarlens}[title={What you've built}]
One more word for telling, explaining and describing at B1.
\end{grammarlens}

\subsection*{Your turn}
\begin{itemize}
\raggedright
  \item \emph{Say it:} \emph{ispolnyát}
  \item \emph{Say it:} \emph{ispolnyát}, once more
  \item \emph{Say it:} a sentence of your own with \emph{ispolnyát}
  \item \emph{Recall:} say the Russian for to conduct, then the Russian for a composer, then say \emph{ispolnyát} again
\end{itemize}

\begin{tcolorbox}[breakable,colback=teal!4,colframe=teal!35!black,title={Before you move on}]
What does \ru{исполнять} mean? (\textquotedblleft{}To perform; to carry out\textquotedblright{}.) Say it once more.
\end{tcolorbox}
\section[\texorpdfstring{\ru{исполнитель} (ispolnítel)}{исполнитель (ispolnítel)}]{\ru{исполнитель} (ispolnítel) — a performer}

\label{lesson:RU-C518-ispolnitel}



\begin{quote}
Before the new one: say the Russian for a composer, then the Russian for to perform.
\end{quote}

\subsection*{You'll want to know: \ru{исполнитель}}
\textbf{\ru{исполнитель}} (\emph{ispolnítel}) — \textquotedblleft{}a performer\textquotedblright{}.

\begin{quote}
\ru{Исполнитель} \ru{вышел} \ru{на} \ru{сцену} \ru{под} \ru{аплодисменты}.
\end{quote}

\begin{quote}
\emph{Ispolnítel výshel na stsénu pod aplodisménty.}
\end{quote}

\begin{quote}
The performer came on stage to applause.
\end{quote}

Masculine; feminine \ru{исполнительница}.

\begin{cousinweb}
\emph{\ru{исполнитель}} is built on \emph{\ru{исполнить}} 'to perform' with \emph{-\ru{тель}}; the root is that of \emph{\ru{полный}} 'full'.
\end{cousinweb}

\begin{grammarlens}[title={What you've built}]
One more word for telling, explaining and describing at B1.
\end{grammarlens}

\subsection*{Your turn}
\begin{itemize}
\raggedright
  \item \emph{Say it:} \emph{ispolnítel}
  \item \emph{Say it:} \emph{ispolnítel}, once more
  \item \emph{Say it:} a sentence of your own with \emph{ispolnítel}
  \item \emph{Recall:} say the Russian for a composer, then the Russian for to perform, then say \emph{ispolnítel} again
\end{itemize}

\begin{tcolorbox}[breakable,colback=teal!4,colframe=teal!35!black,title={Before you move on}]
What does \ru{исполнитель} mean? (\textquotedblleft{}A performer\textquotedblright{}.) Say it once more.
\end{tcolorbox}
\section[\texorpdfstring{\ru{репетиция} (repetítsiya)}{репетиция (repetítsiya)}]{\ru{репетиция} (repetítsiya) — a rehearsal}

\label{lesson:RU-C518-repetitsiya}



\begin{quote}
Before the new one: say the Russian for to perform, then the Russian for a performer.
\end{quote}

\subsection*{You'll want to know: \ru{репетиция}}
\textbf{\ru{репетиция}} (\emph{repetítsiya}) — \textquotedblleft{}a rehearsal\textquotedblright{}.

\begin{quote}
\ru{Репетиция} \ru{начинается} \ru{в} \ru{семь} \ru{часов}.
\end{quote}

\begin{quote}
\emph{Repetítsiya nachináyetsya v sem chasóv.}
\end{quote}

\begin{quote}
The rehearsal starts at seven.
\end{quote}

Feminine.

\begin{cousinweb}
\emph{\ru{репетиция}} was borrowed from Polish \emph{repetycja} (EtymDB), which goes back to Latin \emph{repetītiō} 'a repetition' (Lewis \& Short).
\end{cousinweb}

\begin{grammarlens}[title={What you've built}]
One more word for telling, explaining and describing at B1.
\end{grammarlens}

\subsection*{Your turn}
\begin{itemize}
\raggedright
  \item \emph{Say it:} \emph{repetítsiya}
  \item \emph{Say it:} \emph{repetítsiya}, once more
  \item \emph{Say it:} a sentence of your own with \emph{repetítsiya}
  \item \emph{Recall:} say the Russian for to perform, then the Russian for a performer, then say \emph{repetítsiya} again
\end{itemize}

\begin{tcolorbox}[breakable,colback=teal!4,colframe=teal!35!black,title={Before you move on}]
What does \ru{репетиция} mean? (\textquotedblleft{}A rehearsal\textquotedblright{}.) Say it once more.
\end{tcolorbox}
\section[\texorpdfstring{\ru{репетировать} (repetírovat)}{репетировать (repetírovat)}]{\ru{репетировать} (repetírovat) — to rehearse}

\label{lesson:RU-C518-repetirovat}



\begin{quote}
Before the new one: say the Russian for a performer, then the Russian for a rehearsal.
\end{quote}

\subsection*{You'll want to know: \ru{репетировать}}
\textbf{\ru{репетировать}} (\emph{repetírovat}) — \textquotedblleft{}to rehearse\textquotedblright{}.

\begin{quote}
\ru{Актёры} \ru{репетируют} \ru{новый} \ru{спектакль}.
\end{quote}

\begin{quote}
\emph{Aktyóry repetíruyut nóvyy spektákl.}
\end{quote}

\begin{quote}
The actors are rehearsing a new play.
\end{quote}

Imperfective; present \ru{репетирую}. Perfective \ru{отрепетировать}.

\begin{cousinweb}
\emph{\ru{репетировать}} is built on the stem of \emph{\ru{репетиция}} 'a rehearsal', borrowed from Polish \emph{repetycja} (EtymDB); EtymDB lists the verb without a source.
\end{cousinweb}

\begin{grammarlens}[title={What you've built}]
One more word for telling, explaining and describing at B1.
\end{grammarlens}

\subsection*{Your turn}
\begin{itemize}
\raggedright
  \item \emph{Say it:} \emph{repetírovat}
  \item \emph{Say it:} \emph{repetírovat}, once more
  \item \emph{Say it:} a sentence of your own with \emph{repetírovat}
  \item \emph{Recall:} say the Russian for a performer, then the Russian for a rehearsal, then say \emph{repetírovat} again
\end{itemize}

\begin{tcolorbox}[breakable,colback=teal!4,colframe=teal!35!black,title={Before you move on}]
What does \ru{репетировать} mean? (\textquotedblleft{}To rehearse\textquotedblright{}.) Say it once more.
\end{tcolorbox}
\section[\texorpdfstring{\ru{гастроли} (gastróli)}{гастроли (gastróli)}]{\ru{гастроли} (gastróli) — a tour (by performers)}

\label{lesson:RU-C518-gastroli}



\begin{quote}
Before the new one: say the Russian for a rehearsal, then the Russian for to rehearse.
\end{quote}

\subsection*{You'll want to know: \ru{гастроли}}
\textbf{\ru{гастроли}} (\emph{gastróli}) — \textquotedblleft{}a tour (by performers)\textquotedblright{}.

\begin{quote}
\ru{Летом} \ru{театр} \ru{едет} \ru{на} \ru{гастроли}.
\end{quote}

\begin{quote}
\emph{Létom teátr yédet na gastróli.}
\end{quote}

\begin{quote}
In the summer the theatre company goes on tour.
\end{quote}

Plural. \emph{\ru{Ехать} \ru{на} \ru{гастроли}} — to go on tour.

\begin{cousinweb}
\emph{\ru{гастроли}} is the plural of \emph{\ru{гастроль}}, which was borrowed from German \emph{Gastrolle} 'a guest role' (EtymDB).
\end{cousinweb}

\begin{grammarlens}[title={What you've built}]
One more word for telling, explaining and describing at B1.
\end{grammarlens}

\subsection*{Your turn}
\begin{itemize}
\raggedright
  \item \emph{Say it:} \emph{gastróli}
  \item \emph{Say it:} \emph{gastróli}, once more
  \item \emph{Say it:} a sentence of your own with \emph{gastróli}
  \item \emph{Recall:} say the Russian for a rehearsal, then the Russian for to rehearse, then say \emph{gastróli} again
\end{itemize}

\begin{tcolorbox}[breakable,colback=teal!4,colframe=teal!35!black,title={Before you move on}]
What does \ru{гастроли} mean? (\textquotedblleft{}A tour (by performers)\textquotedblright{}.) Say it once more.
\end{tcolorbox}
